\documentclass[11pt]{article}

\usepackage[margin=1in]{geometry}
\usepackage{amsmath,amssymb,amsthm,mathtools}
\usepackage{booktabs,array}
\usepackage[dvipsnames]{xcolor}
\usepackage{enumitem}
\usepackage{microtype}
\usepackage{fancyhdr}
\usepackage[colorlinks=true,linkcolor=MidnightBlue,urlcolor=MidnightBlue]{hyperref}

\definecolor{ink}{HTML}{17324D}
\definecolor{accent}{HTML}{0F766E}
\definecolor{soft}{HTML}{EEF6F5}
\hypersetup{pdftitle={Finite Abelian Zero-Weight Formula},
  pdfauthor={Computational verification note}}

\newtheorem{theorem}{Theorem}
\newtheorem{proposition}{Proposition}
\newcommand{\CT}{\operatorname{CT}}
\newcommand{\Sym}{\operatorname{Sym}}
\newcommand{\Exp}{\operatorname{Exp}}
\newcommand{\one}{\mathbf 1}

\pagestyle{fancy}
\fancyhf{}
\fancyhead[L]{\small\color{ink}Finite abelian zero-weight formula}
\fancyhead[R]{\small\color{ink}Proof and verification}
\fancyfoot[C]{\thepage}
\setlength{\headheight}{14pt}
\fancypagestyle{plain}{%
  \fancyhf{}%
  \fancyhead[L]{\small\color{ink}Finite abelian zero-weight formula}%
  \fancyhead[R]{\small\color{ink}Proof and verification}%
  \fancyfoot[C]{\thepage}%
}
\fancypagestyle{titlepage}{%
  \fancyhf{}%
  \fancyfoot[C]{\thepage}%
  \renewcommand{\headrulewidth}{0pt}%
}

\title{\color{ink}\bfseries Finite Abelian Representations:\\
The Zero-Weight Formula and a Direct Matrix Verification}
\author{Proof and reproducible verification strategy}
\date{July 17, 2026}

\begin{document}
\maketitle
\thispagestyle{titlepage}

\begin{abstract}
For a finite abelian group acting diagonally, the coefficient of the monomial
symmetric function $m_\tau$ in the sigma moment-generating function is the
number of monomial basis vectors whose total character is trivial.  This note
proves the resulting simultaneous-congruence formula, derives the diagonal
Molien average and finite constant-term form, and records a direct matrix test.
The accompanying marimo notebook constructs the represented group, builds the
full tensor product of symmetric powers, forms the Reynolds projector, and
compares its rank with an exact lattice count and a numerical character average.
\end{abstract}

\section{Setup and the represented matrix group}

Let
\[
 A=C_{m_1}\times\cdots\times C_{m_r},\qquad
 M=\widehat A\cong
 \mathbb Z/m_1\mathbb Z\oplus\cdots\oplus\mathbb Z/m_r\mathbb Z.
\]
Every complex representation of $A$ splits into one-dimensional characters.
Choose a weight basis and write
\[
 V=\bigoplus_{j=1}^{d}\mathbb C_{w_j},\qquad
 w_j=(w_{1j},\ldots,w_{rj})\in M.
\]
The weights are the columns of an integer matrix $W=(w_{\ell j})\in
\mathbb Z^{r\times d}$, understood row by row modulo $m_\ell$.  For
$a=(a_1,\ldots,a_r)$, the corresponding matrix is
\begin{equation}\label{eq:represented-group}
 \rho(a)=\operatorname{diag}_{1\le j\le d}
 \left(\exp\left(2\pi i\sum_{\ell=1}^{r}
 \frac{a_\ell w_{\ell j}}{m_\ell}\right)\right).
\end{equation}
Thus the concrete matrix group is obtained by letting
$0\le a_\ell<m_\ell$ in \eqref{eq:represented-group}.  If $\rho$ has a
kernel, several abstract elements give the same matrix.  These repetitions
must be retained in an average over $A$.

For a partition or finite tuple $\tau=(\tau_1,\ldots,\tau_s)$, put
\[
 T_\tau(V)=\bigotimes_{b=1}^{s}\Sym^{\tau_b}V.
\]
The empty tuple gives $T_\varnothing(V)=\mathbb C$.

\section{The simultaneous-congruence formula}

\begin{theorem}[Zero-weight formula]\label{thm:main}
With the notation above,
\begin{equation}\label{eq:count}
 \dim T_\tau(V)^A=
 \#\left\{(q_{bj})\in\mathbb Z_{\ge0}^{s\times d}:
 \begin{array}{ll}
 \displaystyle\sum_{j=1}^{d}q_{bj}=\tau_b
   & (1\le b\le s),\\[4pt]
 \displaystyle\sum_{b=1}^{s}\sum_{j=1}^{d}
 w_{\ell j}q_{bj}\equiv0\pmod {m_\ell}
   & (1\le\ell\le r)
 \end{array}\right\}.
\end{equation}
Equivalently, for $q_b=(q_{b1},\ldots,q_{bd})^T$ and
$D=\operatorname{diag}(m_1,\ldots,m_r)$,
\begin{equation}\label{eq:lattice}
 \dim T_\tau(V)^A=
 \#\left\{(q_1,\ldots,q_s):
 q_b\in\mathbb Z_{\ge0}^{d},\quad
 \one^Tq_b=\tau_b,\quad
 W(q_1+\cdots+q_s)\in D\mathbb Z^r\right\}.
\end{equation}
\end{theorem}

\begin{proof}
The monomials
\[
 e_1^{q_{b1}}\cdots e_d^{q_{bd}},\qquad
 q_{bj}\ge0,\quad \sum_jq_{bj}=\tau_b,
\]
form a basis of $\Sym^{\tau_b}V$.  Tensoring these bases over $b$ gives a
basis of $T_\tau(V)$ indexed by the arrays in the first line of
\eqref{eq:count}.  Since $e_j$ has character $w_j$, the basis vector indexed
by $(q_{bj})$ has total character
\[
 \sum_{b=1}^{s}\sum_{j=1}^{d}q_{bj}w_j\in M.
\]
It is fixed by every element of $A$ exactly when that character is zero.  In
coordinates this is precisely the second line of \eqref{eq:count}.  Because
the displayed monomial basis is a weight basis, the invariant subspace is
spanned by exactly those basis vectors of weight zero.  Counting them proves
\eqref{eq:count}; collecting the $q_b$ and writing divisibility by each
$m_\ell$ as membership in $D\mathbb Z^r$ gives \eqref{eq:lattice}.
\end{proof}

\section{Relation to the sigma-MGF and Molien averaging}

Let $t_1,t_2,\ldots$ be the alphabet used for the sigma-MGF.  Inserting the
diagonal eigenvalues of \eqref{eq:represented-group} into the generalized
Molien average gives
\begin{equation}\label{eq:mgf}
 \mathbb E_A\!\left[\Exp_\sigma(X_Ah_1)\right]
 =\frac1{|A|}\sum_{a\in A}
 \prod_{i\ge1}\prod_{j=1}^{d}
 \frac{1}{1-t_i\chi_{w_j}(a)}.
\end{equation}
For fixed, distinct alphabet indices, the coefficient of
$t_{i_1}^{\tau_1}\cdots t_{i_s}^{\tau_s}$ in the product for $a$ is
\[
 \prod_{b=1}^{s}h_{\tau_b}
 \bigl(\chi_{w_1}(a),\ldots,\chi_{w_d}(a)\bigr)
 =\operatorname{tr}\bigl(T_\tau(\rho(a))\bigr).
\]
The expression in \eqref{eq:mgf} is symmetric in the $t_i$, so this is also
the coefficient of $m_\tau$.  Therefore
\begin{equation}\label{eq:molien-coefficient}
 [m_\tau]\,\mathbb E_A\!\left[\Exp_\sigma(X_Ah_1)\right]
 =\frac1{|A|}\sum_{a\in A}
 \prod_{b=1}^{s}h_{\tau_b}
 \bigl(\chi_{w_1}(a),\ldots,\chi_{w_d}(a)\bigr)
 =\dim T_\tau(V)^A.
\end{equation}
The final equality follows from the Reynolds projector
\[
 P_A=\frac1{|A|}\sum_{a\in A}T_\tau(\rho(a)).
\]
Indeed, $P_A$ is the projection onto $T_\tau(V)^A$, hence
$\operatorname{tr}P_A=\operatorname{rank}P_A=\dim T_\tau(V)^A$.

There is also a clean constant-term statement.  Work in the group algebra
with $z_\ell^{m_\ell}=1$ and put
$z^{w_j}=\prod_\ell z_\ell^{w_{\ell j}}$.  If $\CT_A$ extracts the coefficient
of the identity character, then
\begin{align}
 \mathbb E_A\!\left[\Exp_\sigma(X_Ah_1)\right]
 &=\CT_A\prod_{i\ge1}\prod_{j=1}^{d}(1-t_i z^{w_j})^{-1},\\
 [m_\tau]\,\mathbb E_A\!\left[\Exp_\sigma(X_Ah_1)\right]
 &=\CT_A\prod_{b=1}^{s}h_{\tau_b}
 (z^{w_1},\ldots,z^{w_d}).
\end{align}
Character orthogonality identifies this operator with the finite Fourier
projector
\[
 \CT_A F=\frac1{|A|}\sum_{a\in A}
 F\bigl(\zeta_{m_1}^{a_1},\ldots,\zeta_{m_r}^{a_r}\bigr).
\]

\section{Verification strategy}

The accompanying notebook uses three routes with different failure modes.

\begin{enumerate}[leftmargin=*,label=\textbf{\arabic*.}]
\item \textbf{Direct represented matrices.}
It constructs all $|A|=\prod_\ell m_\ell$ matrices in
\eqref{eq:represented-group}.  It first checks the identity, unitarity, and
homomorphism laws numerically for every pair of represented elements.  In the
monomial basis, the diagonal of
$\Sym^n(\rho(a))$ is indexed by weak compositions $q$ of $n$ and has entry
$\prod_j\rho(a)_{jj}^{q_j}$.  Kronecker products give the full matrices on
$T_\tau(V)$.  Their average is $P_A$; its numerical rank and trace give the
target quantity directly.  The residual $\lVert P_A^2-P_A\rVert_F$ checks
that the average is a projection.

\item \textbf{Exact zero-weight enumeration.}
It enumerates the same weak compositions using integer arithmetic only and
counts those for which every row of $W(q_1+\cdots+q_s)$ is divisible by the
corresponding $m_\ell$.  This is a literal implementation of
\eqref{eq:lattice}; no complex roots of unity are used.

\item \textbf{Character/Molien average.}
It computes
\[
 \frac1{|A|}\sum_{a\in A}
 \prod_b\operatorname{tr}\bigl(\Sym^{\tau_b}(\rho(a))\bigr)
\]
in floating-point complex arithmetic.  The result must be within $10^{-8}$ of
the exact count.  This route avoids the large Kronecker matrices while still
testing the Fourier projection independently of the modular test.
\end{enumerate}

A case passes when the projector rank, its rounded trace, the exact count, and
the rounded Molien average agree, and both the numerical discrepancy and
idempotence residual are below $10^{-8}$.

\section{Representative results and size choice}

The default case is
\[
 A=C_3\times C_4,\qquad
 W=\begin{pmatrix}1&0&2\\0&1&3\end{pmatrix},\qquad
 d=3,\qquad \tau=(3,2).
\]
Here $|A|=12$ and
\[
\dim T_{(3,2)}(V)=
\binom{3+3-1}{3}\binom{2+3-1}{2}=10\cdot6=60.
\]
\clearpage
\thispagestyle{fancy}
The notebook constructs twelve $60\times60$ target matrices.  It obtains
\[
 \operatorname{rank}P_A=2,\qquad
 \text{zero-weight count}=2,\qquad
 \text{Molien average}=2+3.8\times10^{-15}i,
\]
with maximum displayed numerical discrepancy $8.21\times10^{-15}$.

The automated suite broadens the coverage as follows.
\begin{center}
\small
\begin{tabular}{@{}lrrrrc@{}}
\toprule
Representation & $|A|$ & $d$ & largest $\tau$ & max. target dim. & checks \\
\midrule
$C_2$, trivial plus sign & 2 & 2 & $(3,2)$ & 12 & 5 \\
$C_4$, nonfaithful & 4 & 3 & $(2,1)$ & 18 & 4 \\
$C_2\times C_2$, all characters & 4 & 4 & $(2,1)$ & 40 & 4 \\
$C_3\times C_4$, mixed weights & 12 & 3 & $(3,2)$ & 60 & 5 \\
$C_2\times C_3$, repeated weights & 6 & 4 & $(2,2)$ & 100 & 4 \\
\midrule
\multicolumn{5}{r}{Total} & 22 \\
\bottomrule
\end{tabular}
\end{center}
All 22 checks pass.  The cases include the empty partition, trivial weights,
repeated weights, a nonfaithful representation, one and two cyclic factors,
and tensor products with repeated and unequal degrees.

In general the target dimension is
\begin{equation}\label{eq:target-dim}
 N=\prod_{b=1}^{s}\binom{\tau_b+d-1}{\tau_b}.
\end{equation}
Dense direct projection requires $O(N^2)$ memory and approximately
$O(|A|N^2)$ arithmetic, while literal congruence enumeration visits $N$
monomial basis vectors.  The notebook caps the dense target dimension at
$N=1000$ and uses $N=60$ by default.  For larger dimensions, the Molien trace
average or an exact convolution of residue distributions on
$\prod_\ell\mathbb Z/m_\ell\mathbb Z$ is the appropriate replacement for a
dense projector.

\section{Reproduction}

From the repository root, launch the notebook with
\begin{verbatim}
.venv/bin/marimo edit \
  marimo_notebooks/finite_abelian.py
\end{verbatim}
The controls accept comma-separated moduli and partition parts.  Rows of $W$
are comma-separated and separated from one another by semicolons.  For
example, the default matrix is entered as
\begin{verbatim}
1, 0, 2; 0, 1, 3
\end{verbatim}
The automated table is evaluated whenever the notebook runs.

\end{document}
